%======================================================================
\chapter{Homogeneity and aposyndesis}\label{ch:spread}
%======================================================================

This chapter proves three facts about a homogeneous continuum $X$. The first is a consequence of the
Effros theorem that the formal proof calls \emph{spreading}: a continuum with nonempty interior that
misses a point $x$ can be enlarged to a continuum that still misses $x$ and contains a neighbourhood of
any given one of its points. The second is the Lemma of Jones's 1951 paper: the sets $\Us x$ are never
properly nested. The third is Bing's theorem (1948) that some point of a nondegenerate continuum does
not cut any two open sets from each other; with homogeneity, no point does. These three facts replace
most of the long case analysis of Jones's 1955 paper.

Throughout, $X$ is a compact metric space; from Section~\ref{sec:bing5} on it is also connected.
\emph{Homogeneous} means that for all $x,y\in X$ some homeomorphism $T\colon X\to X$ has $Tx=y$.

\formal{\lean{HPC.Jones.IsHomogeneous}; it is the hypothesis of \lean{HPC.Effros.effros_homogeneous}
(\lean{HPC.Jones.isHomogeneous_iff_isPretransitive}).}

\section{Spreading}

\begin{theorem}[Spreading]\label{thm:spreading}
Let $X$ be a homogeneous compact metric space and $K\subseteq X$ a continuum with nonempty interior. If
$o\in K$ and $x\notin K$, then $X$ is aposyndetic at $o$ with respect to $x$.
\end{theorem}

\begin{proof}
Let $b$ be an interior point of $K$ and choose $\varepsilon_0>0$ with $B(b,\varepsilon_0)\subseteq K$. Let
\[
\varepsilon=\min\bigl(\varepsilon_0,\ \tfrac12\dist(x,K)\bigr)>0 ,
\]
and let $\mathcal S$ be the set of homeomorphisms $h$ of $X$ with $d(h(z),z)<\varepsilon$ for all $z$. Put
\[
Z=\cl{\bigcup_{h\in\mathcal S}\bigl(K\cup h(K)\bigr)} .
\]
\begin{itemize}
\item \emph{$Z$ is a continuum.} Each $K\cup h(K)$ is connected: $h(b)\in B(b,\varepsilon)\subseteq K$, so the
connected sets $K$ and $h(K)$ meet. All these sets contain $o$, so their union is connected, and so is its
closure. $Z$ is closed in the compact space $X$.
\item \emph{$x\notin Z$.} Every point of $h(K)$ is within $\varepsilon$ of $K$, so every point of $Z$ is at
distance $\le\varepsilon\le\frac12\dist(x,K)$ from $K$.
\item \emph{$Z$ contains a neighbourhood of $o$.} By the Effros theorem (\vthree{Theorem~7.15}) there is
$\delta>0$ such that every $y$ with $d(o,y)<\delta$ is $h(o)$ for some $h\in\mathcal S$; then $y\in h(K)\subseteq Z$. So
$B(o,\delta)\subseteq Z$.\qedhere
\end{itemize}
\end{proof}

\formal{\lean{HPC.Jones.aposAt_of_isContinuum_interior}, from \lean{HPC.Effros.effros_homogeneous}.}

\begin{corollary}[Non-aposyndesis is transitive]\label{cor:trans}
Let $X$ be homogeneous. If $X$ is not aposyndetic at $z$ with respect to $y$, nor at $y$ with respect to
$x$, then it is not aposyndetic at $z$ with respect to $x$. In terms of the sets of
Definition~\ref{def:apos}: if $z\in\Ls y$ and $y\in\Ls x$, then $z\in\Ls x$.
\end{corollary}

\begin{proof}
Suppose $K$ is a continuum missing $x$ that contains an open $V\ni z$. Then $y\in K$, for otherwise $(K,V)$
would make $X$ aposyndetic at $z$ with respect to $y$. Now $K$ has nonempty interior, contains $y$ and misses
$x$; by Theorem~\ref{thm:spreading}, $X$ is aposyndetic at $y$ with respect to $x$, a contradiction.
\end{proof}

\formal{\lean{HPC.Jones.not_aposAt_trans}.}

\section{The Lemma of Jones 1951}

\begin{lemma}\label{lem:open-subset-U}
For every closed $K\subseteq X$, the set $\set{w\in X : K\subseteq\Us w}$ is open.
\end{lemma}

\begin{proof}
Let $K\subseteq\Us w$, that is, $X$ is aposyndetic at every point of $K$ with respect to $w$. By
Lemma~\ref{lem:heine-borel}(a), finitely many continua $H_i$ missing $w$ have interiors covering $K$. For $w'$
in the open set $X\setminus\bigcup_iH_i$ the same continua show $K\subseteq\Us{w'}$.
\end{proof}

\formal{\lean{HPC.Jones.isOpen_setOf_subset_aposSet}.}

\begin{theorem}[Jones 1951, Lemma]\label{thm:jones51-lemma}
If $X$ is homogeneous, then $\Us y$ is not a proper subset of $\Us x$, for any $x,y\in X$.
\end{theorem}

\begin{proof}
Let $B_0,B_1,\dots$ be a countable basis of $X$ and define
\[
g(w)=\sum_{k\ge0}2^{-k}\,\chi_k(w),\qquad \chi_k(w)=\begin{cases}1&\text{if }\cl{B_k}\subseteq\Us w,\\ 0&\text{otherwise.}\end{cases}
\]
By Lemma~\ref{lem:open-subset-U} each $\chi_k$ is the indicator function of an open set, hence lower
semicontinuous, and so is the sum $g$. A lower semicontinuous function on a compact space attains its
minimum, say at $w_0$.

If $\Us a\subsetneq\Us b$, then $g(a)<g(b)$: every term of $g(a)$ is at most the corresponding term of $g(b)$;
and for $z\in\Us b\setminus\Us a$ there is $k$ with $z\in B_k\subseteq\cl{B_k}\subseteq\Us b$ (the set $\Us b$ is open,
Theorem~\ref{thm:jones48-1}), and then $\chi_k(b)=1$, $\chi_k(a)=0$ because $z\in\cl{B_k}\setminus\Us a$.

Now suppose $\Us y\subsetneq\Us x$ and let $T$ be a homeomorphism with $Tx=w_0$. By
Lemma~\ref{lem:apos-basic}(b), $\Us{Ty}=T(\Us y)\subsetneq T(\Us x)=\Us{w_0}$, so $g(Ty)<g(w_0)$, against the
minimality of $w_0$.
\end{proof}

\formal{\lean{HPC.Jones.jones51_lemma}, with the function $g$ from
\lean{HPC.Jones.exists_lsc_strictMono_aposSet}. Jones's proof runs a transfinite sequence of points
$x_0,x_1,\dots$ through the countable ordinals with $\Us{x_i}$ strictly decreasing, and contradicts second
countability; the lower semicontinuous function replaces the transfinite recursion.}

\section{Bing's theorem on non-cutting points}\label{sec:bing5}

From now on $X$ is a nondegenerate continuum. Bing's paper \emph{Some characterizations of arcs and simple
closed curves} (1948) contains the following theorem, which Jones's 1955 paper quotes as
``Corollary~2 of [Bing]''.

\begin{theorem}[Bing 1948, Theorem~5, point form]\label{thm:bing5}
For every $r\in X$ there is a point $p\ne r$ such that the union of the subcontinua of $X$ which contain $r$
and miss $p$ is dense in $X$.
\end{theorem}

The proof needs a definition and three lemmas. Call a subcontinuum $N$ of $X$ \emph{tight} if there are no
proper subcontinuum $K$ of $X$ and open set $O$ with $N\subseteq O\subseteq K$; that is, $N$ lies in the interior of no
proper subcontinuum.

\begin{lemma}\label{lem:tight-conn}
Let $N\subseteq E$ be subcontinua of $X$ with $N$ tight. Then $\cl{X\setminus E}\cup N$ is connected.
\end{lemma}

\begin{proof}
Let $G=\cl{X\setminus E}\cup N$ and suppose $G=A\cup B$ with $A$, $B$ closed, disjoint and nonempty; as $N$ is
connected we may take $N\subseteq A$. Choose disjoint open sets $u\supseteq A$, $v\supseteq B$ and put $Q=A\cup E$, a closed set.

\emph{$Q$ is connected.} Otherwise $Q=a\cup b$ with $a$, $b$ closed, disjoint and nonempty, and $E\subseteq a$; then
$b\subseteq A\setminus E$. Every point of $X\setminus Q$ lies in $X\setminus E\subseteq G$ but not in $A$, hence in $B$. So
$X=b\cup(a\cup B)$, where $b$ and $a\cup B$ are closed, nonempty and disjoint (because $b\subseteq A$). This
contradicts the connectedness of $X$.

\emph{$Q\ne X$.} $B$ is nonempty and contained in $\cl{X\setminus E}$ (as $B\cap N=\emptyset$), so the open set $v\supseteq B$
contains a point $z\notin E$; and $z\notin A$ since $A\subseteq u$. So $z\notin Q$.

\emph{$u\subseteq Q$.} A point of $u$ outside $E$ lies in $G=A\cup B$ but not in $B\subseteq v$, so it lies in $A$.

Thus $Q$ is a proper subcontinuum containing the open set $u\supseteq N$, against the tightness of $N$.
\end{proof}

\formal{\lean{HPC.Plane.Arcs48.isPreconnected_closure_diff_union} (Bing: ``$F_i+N$ is a continuum'').}

\begin{lemma}\label{lem:tight-piece}
Let $N$ be a tight subcontinuum and $D$ a nonempty open set with $D\cap N=\emptyset$. There is a closed set $Z\subseteq X$
with empty interior such that for every $p\in X\setminus(N\cup Z)$ some continuum $G\supseteq N$ misses $p$ and meets $\cl
D$.
\end{lemma}

\begin{proof}
Fix $n_0\in N$. The closed set $X\setminus D$ is proper and contains $n_0$. By boundary bumping
(\vthree{Theorem~1.10}) the component $K$ of $n_0$ in $X\setminus D$ meets $\cl D$. Let $E=K\cup N$, a continuum
contained in $X\setminus D$, and $Z=E\cap\cl{X\setminus E}$. A nonempty open set inside $Z$ would lie in $E$ and meet
$X\setminus E$ (being inside $\cl{X\setminus E}$); so $Z$ has empty interior. Let $p\notin N\cup Z$. If $p\notin E$,
take $G=E$. If $p\in E$, then $p\notin\cl{X\setminus E}$; take $G=\cl{X\setminus E}\cup N$, a continuum by
Lemma~\ref{lem:tight-conn}, which contains $D\subseteq X\setminus E$ and misses $p$.
\end{proof}

\formal{\lean{HPC.Plane.Arcs48.exists_nowhereDense_piece}.}

\begin{lemma}[The tight case]\label{lem:tight-case}
If $N$ is a tight proper subcontinuum, there is $p\in X\setminus N$ such that the union of the continua containing $N$
and missing $p$ is dense in $X$.
\end{lemma}

\begin{proof}
Let $b_0,b_1,\dots$ be a countable basis. For each $k$ with $b_k\ne\emptyset$ and $\cl{b_k}\cap N=\emptyset$ let $Z_k$ be the
set of Lemma~\ref{lem:tight-piece} for $D=b_k$. The sets $Z_k$ are closed with empty interior and $X\setminus N$ is a
nonempty open set, so by the Baire category theorem some $p\in X\setminus N$ lies in no $Z_k$. Let $x\in X$ and $W\ni x$
open. If $x\in N$, the continuum $N$ itself meets $W$. Otherwise choose $k$ with $x\in b_k\subseteq\cl{b_k}\subseteq W\setminus N$; the
continuum $G$ given by Lemma~\ref{lem:tight-piece} contains $N$, misses $p$ and meets $\cl{b_k}\subseteq W$.
\end{proof}

\formal{\lean{HPC.Plane.Arcs48.exists_point_of_tight}.}

\begin{lemma}[Bing's sequence]\label{lem:bing-seq}
Let $N$ be a nonempty proper subcontinuum. Either there is $p\in X\setminus N$ such that the continua containing $N$ and
missing $p$ have dense union, or some tight proper subcontinuum contains $N$.
\end{lemma}

\begin{proof}
Let $b_0,b_1,\dots$ be a countable basis. Put $N_0=N$, and let $N_{j+1}$ be a proper subcontinuum $K$ such that some
open $O$ has $N_j\cup b_j\subseteq O\subseteq K$, if such a $K$ exists; otherwise $N_{j+1}=N_j$. The $N_j$ form an increasing
sequence of proper subcontinua containing $N$; let $U=\bigcup_jN_j$, a connected set.

\emph{If $U$ is dense, then $U\ne X$.} Suppose $U=X$, and let $W_n$ be the interior of $N_n$; the $W_n$ increase. If
$x\in N_i$, some $j\ge i$ is a step at which $K$ exists (otherwise $N_j=N_i$ for all $j\ge i$ and $N_i=U=X$), and then
$x\in N_j\subseteq O\subseteq N_{j+1}$, so $x\in W_{j+1}$. By compactness $X=W_n$ for some $n$, so $N_n=X$, which is false. Any $p\notin
U$ then satisfies the first alternative, with the continua $N_n$.

\emph{If $U$ is not dense, $N'=\cl U$ is tight.} $N'$ is a proper subcontinuum containing $N$. Let $K$ be a proper
subcontinuum and $O$ open with $N'\subseteq O\subseteq K$. Some $q\in O$ is not in $N'$: otherwise $N'=O$ would be open and closed,
nonempty and not all of $X$. Choose $j$ with $q\in b_j\subseteq O\setminus N'$. Then $N_j\cup b_j\subseteq O\subseteq K$, so at step $j$ a
continuum $N_{j+1}$ was chosen with $b_j\subseteq N_{j+1}$, and $q\in U\subseteq N'$, a contradiction.
\end{proof}

\formal{\lean{HPC.Plane.Arcs48.tight_or_point}, \lean{HPC.Plane.Arcs48.exists_point_dense}.}

\begin{proof}[Proof of Theorem~\ref{thm:bing5}]
$N=\set r$ is a proper subcontinuum. In the first case of Lemma~\ref{lem:bing-seq} we are done; in the second,
apply Lemma~\ref{lem:tight-case} to the tight continuum $N'\ni r$.
\end{proof}

\formal{\lean{HPC.Plane.arcs_theorem5_point}; the general form for a proper subset $R$ in place of $\set r$ is
\lean{HPC.Plane.arcs_theorem5}.}

\begin{corollary}[Bing 1948, Corollary~2]\label{cor:bing2}
A nondegenerate continuum $X$ contains two points $p\ne q$ with the following property: for any nonempty open sets
$D,E\subseteq X$ there is a continuum that meets $D$ and $E$ and misses $p$ (respectively $q$).
\end{corollary}

\begin{proof}
Let $r\in X$ and let $p$ be given by Theorem~\ref{thm:bing5}. Continua $K_1,K_2$ containing $r$ and missing $p$ meet $D$,
respectively $E$; their union is the required continuum. Then let $q\ne p$ be given by Theorem~\ref{thm:bing5} for
$r=p$.
\end{proof}

\formal{\lean{HPC.Plane.arcs_corollary2}, with Bing's relation ``$R$ cuts $A$ from $B$ in $M$'' as
\lean{HPC.Plane.Cuts}.}

\section{Consequences for homogeneous continua}

\begin{proposition}[Jones 1955, claims (3) and (5)]\label{prop:claim35}
Let $X$ be a nondegenerate homogeneous continuum and $p\in X$.
\begin{enumerate}[label=(\alph*)]
\item For any nonempty open sets $D,E$ there is a continuum meeting $D$ and $E$ and missing $p$.
\item Any two points $a,b\in\Us p$ lie in a continuum missing $p$.
\end{enumerate}
\end{proposition}

\begin{proof}
(a) Let $p_0$ be a point given by Corollary~\ref{cor:bing2} and $T$ a homeomorphism with $Tp_0=p$. A continuum
meeting $T^{-1}(D)$ and $T^{-1}(E)$ and missing $p_0$ is carried by $T$ to the required continuum.

(b) Let $A\supseteq V\ni a$ and $B\supseteq W\ni b$ be continua missing $p$ containing open sets $V$, $W$. By (a) a continuum $K$
missing $p$ meets $V$ and $W$. Then $A\cup K\cup B$ is a continuum containing $a$, $b$ and missing $p$.
\end{proof}

\formal{\lean{HPC.Jones.exists_continuum_meets_of_open}, \lean{HPC.Jones.jones55_claim3},
\lean{HPC.Jones.exists_continuum_of_mem_aposSet}, \lean{HPC.Jones.jones55_claim5}.}

\section*{Exercises}

\begin{exercise}
Show that Theorem~\ref{thm:spreading} fails for the arc $[0,1]$ (which is not homogeneous): find a continuum $K$ with
nonempty interior, $o\in K$ and $x\notin K$ such that every continuum containing a neighbourhood of $o$ contains $x$.
\end{exercise}

\begin{exercise}
In the arc $[0,1]$, which points $p$ satisfy the conclusion of Corollary~\ref{cor:bing2}? Show that in the circle every
point does.
\end{exercise}

\begin{exercise}
Prove that a lower semicontinuous function $g\colon X\to[0,\infty]$ on a compact space attains its minimum.
\end{exercise}
